% 三种融合位置的结构示意；宽约150mm，跨栏使用。只画模块与数据流，不表示层数。
\begin{tikzpicture}[x=1cm,y=1cm,font=\figurefont\small,
  block/.style={draw=FigureModel,fill=FigureModel!10,line width=0.8pt,
    minimum height=8mm,minimum width=20mm,align=center,rounded corners=1.8pt},
  input/.style={block,draw=FigureInput,fill=FigureInput!10,minimum width=15mm},
  output/.style={block,draw=FigureOutput,fill=FigureOutput!10,minimum width=14mm},
  arr/.style={-{Stealth[length=2mm]},BookInk,line width=0.8pt}]
  % 早期融合
  \node[anchor=west,font=\figurefont\small\bfseries] at (0,8.2) {早期融合：先拼接特征，再由一个网络预测};
  \node[input] (ea) at (1.0,7.3) {音频$\bm z_{\mathrm a}$};
  \node[input] (ev) at (1.0,6.2) {视觉$\bm z_{\mathrm v}$};
  \node[block] (ec) at (4.6,6.75) {拼接\\$[\bm z_{\mathrm a};\bm z_{\mathrm v}]$};
  \node[block] (ef) at (8.4,6.75) {共享网络$f$};
  \node[output] (ey) at (12.0,6.75) {预测$\hat y$};
  \draw[arr] (ea) -- (ec);
  \draw[arr] (ev) -- (ec);
  \draw[arr] (ec) -- (ef);
  \draw[arr] (ef) -- (ey);
  % 晚期融合
  \node[anchor=west,font=\figurefont\small\bfseries] at (0,5.1) {晚期融合：各模态先给出决策，再汇总};
  \node[input] (la) at (1.0,4.2) {音频$\bm z_{\mathrm a}$};
  \node[input] (lv) at (1.0,3.1) {视觉$\bm z_{\mathrm v}$};
  \node[block] (lfa) at (4.6,4.2) {音频网络$f_{\mathrm a}$};
  \node[block] (lfv) at (4.6,3.1) {视觉网络$f_{\mathrm v}$};
  \node[block] (lc) at (8.4,3.65) {决策汇总\\加权或相乘};
  \node[output] (ly) at (12.0,3.65) {预测$\hat y$};
  \draw[arr] (la) -- (lfa);
  \draw[arr] (lv) -- (lfv);
  \draw[arr] (lfa) -- node[above,sloped] {$p_{\mathrm a}$} (lc);
  \draw[arr] (lfv) -- node[below,sloped] {$p_{\mathrm v}$} (lc);
  \draw[arr] (lc) -- (ly);
  % 中间融合
  \node[anchor=west,font=\figurefont\small\bfseries] at (0,2.0) {中间融合：在中间表示上建立显式交互};
  \node[input] (ma) at (1.0,1.1) {音频$\bm z_{\mathrm a}$};
  \node[input] (mv) at (1.0,0.0) {视觉$\bm z_{\mathrm v}$};
  \node[block] (mi) at (4.6,0.55) {交互模块\\双线性／FiLM／交叉注意力};
  \node[block] (mg) at (8.4,0.55) {后续网络$g$};
  \node[output] (my) at (12.0,0.55) {预测$\hat y$};
  \draw[arr] (ma) -- (mi);
  \draw[arr] (mv) -- (mi);
  \draw[arr] (mi) -- (mg);
  \draw[arr] (mg) -- (my);
\end{tikzpicture}
